\section{Boundary preparation and independent supports}
\label{sec:preparation}

We now work in the setting furnished by Proposition~\ref{prop:smoothness}:
$V\subset\C^N$ is irreducible of complex dimension $n\geq1$, and
$U\subset V_{\mathrm{reg}}$ is a domain.  The purpose of this section is
to choose a boundary point with two simultaneous properties.  Its normal
fibers have limits uniform in the tangential parameters, and it admits
enough global holomorphic supports to control all complex-normal directions.
Neither property is asserted for an arbitrary supported boundary point.

We use the elementary semialgebraic facts of cell decomposition, selection,
and generic real analyticity.  More precisely, a semialgebraic set admits
a finite decomposition into connected real-analytic cells with algebraic
analytic parametrizations; a choice in nonempty semialgebraic fibers can
be made semialgebraically; and a semialgebraic function on an open set is
algebraic analytic off a lower-dimensional set.  For projection and
real-analytic cell decomposition, see
\cite[Sections~1.1--1.2]{Coste2005}.  These facts apply equally
to descriptions using real quantifiers.  Lower-dimensional exceptional
sets have Lebesgue measure zero.

We first separate boundary cells that admit normal graph coordinates
from those contained in a proper complex algebraic subset.  The latter
can be excluded from a maximizing construction by a polynomial factor.

\begin{lemma}[Complex rank of a boundary cell]
\label{lem:rank-dichotomy}
Let $S$ be a connected real-analytic semialgebraic cell in
$V_{\mathrm{reg}}$.  Either
\begin{equation}
 T_pS+iT_pS=T_pV
 \label{eq:generic-cell}
\end{equation}
outside an intrinsic null set of $S$, or $S$ is contained in a proper
complex algebraic subset of $V$.
\end{lemma}

\begin{proof}
Put $d=\dim_{\R}S$, and write the ambient complex coordinate functions
of an algebraic analytic parametrization of $S$ as
$q_1(s),\ldots,q_N(s)$, where $s=(s_1,\ldots,s_d)$.  Complexify the
parameters.  These germs belong to a common finite algebraic extension
of $\C(s_1,\ldots,s_d)$.  In characteristic zero, the dimension of the
span of their differentials is
\[
 \operatorname{trdeg}_{\C}\C(q_1,\ldots,q_N).
\]
Indeed, a transcendence basis among the $q_j$ has independent
differentials, as seen by extending it to a transcendence basis of the
ambient field; all remaining $q_j$ are separably algebraic over it and
add no independent differentials.  This number is also the generic
complex rank of the complexified parametrization and the dimension of
its complex Zariski closure.  The latter lies in $V$.

If the rank is less than $n$, a polynomial vanishing on that closure
but not identically on $V$ vanishes on an open patch of $S$, hence on
all of $S$ by real-analytic continuation.  If the rank is $n$, its
failure locus is defined locally by nontrivial analytic minors and has
intrinsic measure zero.  The complex image of the differential is
exactly $T_pS+iT_pS$, which proves the alternative.
\end{proof}

Decompose $\partial U\cap V_{\mathrm{reg}}$ into finitely many such
cells.  Every cell has positive real codimension, since the boundary of
an open set has empty interior.  On a cell satisfying
Equation~\eqref{eq:generic-cell}, write
\[
 k=2n-\dim_{\R}S,\qquad m=n-k,
 \qquad 1\leq k\leq n.
\]
At a point satisfying Equation~\eqref{eq:generic-cell}, choose
semialgebraic holomorphic coordinates, obtained locally from a linear
projection of $V$ followed by a complex-linear change, in which a patch
of $S$ is
\begin{equation}
 \operatorname{Im}w_0=\phi(s),\qquad
 s=(z_0,\operatorname{Re}w_0)\in B,
 \qquad (z_0,w_0)\in\C^m\times\C^k.
 \label{eq:boundary-graph}
\end{equation}
Here $B$ is an open real parameter box and $\phi$ is real analytic and
semialgebraic.  This follows from the real implicit function theorem
after putting the tangent space in the form
$\{\operatorname{Im}w_0=0\}$.  Shrinking the patch ensures that its
intersection with each fixed-$z_0$ slice is totally real in $\C^k$.

Choose $\rho>0$ so that the points with $s\in B$ and
$|v|<\rho$ below stay in the coordinate patch, and put
\begin{equation}
 \begin{split}
 v&=\operatorname{Im}w_0-\phi(s),\qquad
 s=(z_0,x_0),\quad x_0=\operatorname{Re}w_0,\\
 A_s&=\{v\in\R^k:|v|<\rho,\
       (z_0,x_0+i(\phi(s)+v))\in U\}.
 \end{split}
 \label{eq:offset-fibers}
\end{equation}
Thus the $A_s$ are open semialgebraic fibers, and $0\notin A_s$.
The cutoff $\rho$ will disappear under every bounded-coordinate scaling
as $t\downarrow0$.

We need two properties of these fibers at the same parameter: analytic
families of offsets in $A_s$ approaching zero, and distance functions
for the complements of $t^{-1}A_s$ that stabilize as $s$ moves and
$t$ tends to zero.  The joint limit matters because later coordinate
changes move the tangential parameter at the same time as the normal
scale.  The following parameter version of ramification supplies both
properties outside a null set.

\begin{lemma}[Ramification with parameters]
\label{lem:ramified-parameters}
Let $B\subset\R^d$ be a semialgebraic open set and let
$b:B\times(0,t_0)\to\R$ be a bounded semialgebraic function.  There is a
null set $E\subset B$ such that, for every $s_*\in B\setminus E$, there
are a neighborhood $B_*$ of $s_*$, an integer $l\geq1$, a number
$\eps>0$, and a real-analytic function $\widetilde b$ on
$B_*\times(-\eps,\eps)$ satisfying
\[
 b(s,u^l)=\widetilde b(s,u)
 \qquad(s\in B_*,\ 0<u<\eps).
\]
In particular, $b(s,t)$ has a limit as $(s,t)\to(s_*,0)$ with $t>0$,
without any restriction on the relative rates of approach.

Finitely many such functions can use a common ramification and
neighborhood.  For countably many functions there is a common full-measure
set of good points, with the conclusion applied separately to each
function; no common ramification or neighborhood is asserted.
\end{lemma}

\begin{proof}
Use cylindrical decomposition with $s$ preceding $t$.  Away from a
lower-dimensional subset of $B$, the region immediately above $t=0$
lies in a cell on which $b$ is an algebraic analytic branch.  Locally at
such an $s_*$, shrink to a rectangle $B_*\times(0,\eps_0)$ on which
this holds.  Choose a square-free polynomial equation
\[
 R(s,t,b(s,t))=0
\]
with polynomial coefficients in $(s,t)$ and nonzero leading coefficient
and discriminant in the last variable.  Remove the zero sets of the
first nonzero $t$-coefficients of these latter two polynomials.  After
shrinking again, each is a power of $t$ times a nonvanishing holomorphic
function on a complex $s$-polydisc times a complex $t$-disc.

For $t\ne0$ the roots therefore form a finite unramified covering of
this product with the $t$-origin deleted.  The $s$-polydisc is simply
connected, so its monodromy is generated by one finite permutation around
$t=0$.  The substitution $t=u^l$ kills this permutation for some $l$.
The chosen branch becomes holomorphic for $u\ne0$.  The polynomial root
bound, applied after division by the leading coefficient, bounds this
branch by $C|u|^{-M}$ on a smaller complex polydisc, for some finite $M$.
Thus it has at worst a pole at $u=0$, whose negative Laurent coefficients
are holomorphic in $s$.  Boundedness of the original function for real
$s$ and positive $u$ makes those coefficients vanish for every real $s$
in the rectangle.  They vanish identically by holomorphic uniqueness.
The branch consequently extends holomorphically through $u=0$, giving
the stated real-analytic extension.

There are only finitely many decomposition branches and exceptional
polynomial zero sets.  This proves the assertion outside a null set.
For finitely many functions take a common multiple of their exponents
and shrink the neighborhoods.  For countably many, discard the union
of their exceptional null sets; the individual conclusions suffice.
\end{proof}

\begin{proposition}[Generic boundary data]
\label{prop:boundary-data}
In the coordinates of Equation~\eqref{eq:boundary-graph}, almost every
$s_*\in B$ has the following two properties.
\begin{enumerate}[label=\textup{(\roman*)}]
\item There are a neighborhood $B_*\subset B$ of $s_*$, a number
$\eps>0$, and a real-analytic map
$\psi:B_*\times(-\eps,\eps)\to\R^k$ such that
\[
 \psi(s,0)=0,\qquad
 \psi(s,u)\in A_s\quad(s\in B_*,\ 0<u<\eps).
\]
\item Set
\[
 E(s,t)=\R^k\setminus t^{-1}A_s,
 \qquad d(s,t;v)=\min\{1,\operatorname{dist}(v,E(s,t))\}.
\]
There is a $1$-Lipschitz function $d_*:\R^k\to[0,1]$ such that,
for every compact $L\subset\R^k$,
\begin{equation}
 \lim_{\substack{s\to s_*\\t\downarrow0}}
       \sup_{v\in L}|d(s,t;v)-d_*(v)|=0.
 \label{eq:joint-offset-limit}
\end{equation}
The set $F=\{v:d_*(v)=0\}$ is a closed cone, invariant under every
positive real dilation, and $0\in F$.
\end{enumerate}
In particular, for every sequence $s_j\to s_*$, $t_j\downarrow0$,
and $v_j\to v$, one has $d(s_j,t_j;v_j)\to d_*(v)$.  If $v\in F$,
then
\[
 \operatorname{dist}(v_j,E(s_j,t_j))\longrightarrow0.
\]
If $v\notin F$, some fixed ball about $v$ is contained in
$t_j^{-1}A_{s_j}$ for all sufficiently large $j$.
\end{proposition}

\begin{proof}
The function
\[
 a(s)=\min\{1,\operatorname{dist}(0,A_s)\},
\]
with value $1$ for an empty fiber, is semialgebraic and continuous on
an open full-measure subset $B'\subset B$.  It is zero there.  In fact,
if $a(s_0)>0$ at a continuity point, it is bounded below near $s_0$;
but the boundary point with parameters $(s_0,0)$ is a limit of points
of $U$, whose parameters $(s_j,v_j)$ satisfy $s_j\to s_0$ and
$v_j\to0$.  This is a contradiction.

On $B'\times(0,t_0)$ semialgebraic selection therefore supplies
$v(s,t)\in A_s$ with $|v(s,t)|<t$.  Apply
Lemma~\ref{lem:ramified-parameters} to its finitely many coordinates.
At almost every $s_*$ this gives a jointly analytic
$\psi(s,u)=v(s,u^l)$.  The bound $|\psi(s,u)|<u^l$ proves
$\psi(s,0)=0$ for every $s$ in its neighborhood.  This proves (i).

For (ii), apply Lemma~\ref{lem:ramified-parameters} separately to
$d(s,t;v)$ for every $v\in\mathbb{Q}^k$.  Their exceptional sets have
a null union.  At a remaining $s_*$, each rational test has a limit
along every joint approach $(s,t)\to(s_*,0)$, $t>0$.  For every $s,t$,
\[
 |d(s,t;v)-d(s,t;v')|\leq |v-v'|.
\]
The rational limits therefore extend uniquely to a $1$-Lipschitz
function $d_*$ on $\R^k$.  A finite rational net in a compact set,
together with this inequality, proves
Equation~\eqref{eq:joint-offset-limit}.  Notice that this argument
uses only finitely many test neighborhoods at a time; it does not
require a common neighborhood for all rational tests.

The zero set $F$ is closed and contains $0$, since $0\in E(s,t)$.
For $c>0$ the exact identity
\[
 E(s,ct)=c^{-1}E(s,t)
\]
shows, on taking $s=s_*$ and $t\downarrow0$, that the distance from
$v$ to the left-hand side tends to zero if and only if the distance
from $cv$ to $E(s_*,t)$ does.  Truncation at $1$ does not affect
vanishing of a limit.  Consequently $v\in F$ if and only if $cv\in F$.

The assertion for moving vectors follows from the same Lipschitz
inequality.  If $d_*(v)=0$, convergence of the truncated distance to
zero is convergence of the untruncated distance to zero.  If
$d_*(v)>0$, the distance of $v$ from $E(s_j,t_j)$ is eventually at
least $d_*(v)/2$, so any smaller fixed ball is in the scaled interior.
\end{proof}

Using countably many patches of the form
Equation~\eqref{eq:boundary-graph}, we obtain an intrinsic full-measure
set of points on every cell in the first alternative of
Lemma~\ref{lem:rank-dichotomy}, each with the data of
Proposition~\ref{prop:boundary-data}.  Call these points \emph{good}.
The cone $F$ at a good point might still equal all of $\R^k$;
Proposition~\ref{prop:noncollapse} will rule this out.  We first choose
a good point carrying independent supports.

\begin{proposition}[A good supported point]
\label{prop:supported-point}
There are a boundary cell $S$ in the first alternative of
Lemma~\ref{lem:rank-dichotomy}, a good point $p\in S$, and
$k=2n-\dim_{\R}S$ entire functions $h_1,\ldots,h_k$ on $\C^N$
such that
\begin{equation}
 h_j(p)=1,\qquad
 |h_j(q)|\leq\exp(-|q-p|^2)\quad(q\in\overline U),
 \qquad 1\leq j\leq k.
 \label{eq:peak}
\end{equation}
The real covectors
\[
 \left.d\log|h_j|\right|_{T_pV}
\]
are linearly independent and annihilate $T_pS$.  Consequently, in
coordinates $(z,w)\in\C^{n-k}\times\C^k$ centered at $p$ with
$T_pS=\{\operatorname{Im}w=0\}$, their holomorphic logarithmic
differentials have the form
\[
 d\log h_j(p)=i\beta_j\,dw,
\]
where the $\beta_j$ are linearly independent real rows.
\end{proposition}

\begin{proof}
We vary the real linear term in the weighted maximum used in the
smoothness proof.  We seek one good point at which the maximum is
attained for a positive-measure set of linear terms, measured within
the set of all terms making that point critical on its boundary cell.
The corresponding conormals will provide $k$ independent supports.

Choose a polynomial $P$ vanishing on $V_{\mathrm{sing}}$ and on the
proper complex algebraic subsets containing all cells in the second
alternative of Lemma~\ref{lem:rank-dichotomy}, but not identically on
$V$.  Such a polynomial exists because the union is finite and $V$ is
irreducible.  For $a\in\C^N$, maximize on the compact set
$\overline U$ the upper-semicontinuous function
\[
 \Phi_a(q)=\log|P(q)|+|q|^2+\operatorname{Re}(a\cdot q),
\]
with value $-\infty$ where $P=0$.  Its maximum $M(a)$ is finite,
because $P$ is not identically zero on the open set $U$.  Every
maximizer has $P\ne0$, lies in $V_{\mathrm{reg}}$, and belongs to
$\partial U$: on the smooth locus with $P\ne0$, strict
plurisubharmonicity forbids an interior maximum.  Thus every maximizer
lies on a retained boundary cell.

Fix one retained cell $S$, of real dimension $d$, and restrict to
$S^\circ=S\cap\{P\ne0\}$.  Its critical-pair space is
\[
 \mathcal E_S=
 \{(p,a)\in S^\circ\times\C^N:
                  d(\Phi_a|_S)_p=0\}.
\]
Ambient real linear functionals restrict onto $T_p^*S$, so these are
$d$ independent real affine equations on the $2N$ real coordinates
of $a$.  Therefore $\mathcal E_S\to S^\circ$ is a smooth affine
bundle with fiber dimension $2N-d$ and total dimension $2N$.

The subset $\mathcal M_S\subset\mathcal E_S$ of actual maximizers
is Borel.  Indeed, $M(a)$ is Lipschitz in $a$, since
\[
 |M(a)-M(a')|
 \leq\bigl(\sup_{q\in\overline U}|q|\bigr)|a-a'|,
\]
and $\mathcal M_S$ is defined by the continuous equality
$\Phi_a(p)=M(a)$ on $\mathcal E_S$.  Projection
$\mathcal E_S\to\C^N$ is a smooth map between manifolds of the
same real dimension and sends null sets to null sets: in countably
many relatively compact coordinate charts it is Lipschitz.  Its
restrictions to the finitely many sets $\mathcal M_S$ cover all of
$\C^N$.  Hence at least one $\mathcal M_S$ has positive measure in
$\mathcal E_S$.

In local bundle coordinates, the points lying over the nongood null
subset of $S$ form a null set, by Fubini.  Applying Fubini again to
$\mathcal M_S$, in countably many bounded bundle charts, gives a
good point $p$ for which the set of actual maximizing parameters in
the critical fiber has positive $(2N-d)$-dimensional measure.  Fix
this $p$ and write that subset as $\mathcal A_p$.

Each $a\in\mathcal A_p$ supplies the entire function
\[
 h_a(q)=\frac{P(q)}{P(p)}
       \exp\bigl((2\overline p+a)\cdot(q-p)\bigr).
\]
The inequality $\Phi_a(q)\leq\Phi_a(p)$ gives
\[
 \log|h_a(q)|
 \leq |p|^2-|q|^2
       +2\operatorname{Re}\bigl(\overline p\cdot(q-p)\bigr)
 =-|q-p|^2,
\]
with the desired modulus inequality also at $P(q)=0$.

The affine map from the full critical fiber to the conormal space
\[
 a\longmapsto
 \left.d\log|h_a|\right|_{T_pV}
 \in N_p^*S
 :=\{\lambda\in T_p^*V:\lambda|_{T_pS}=0\}
\]
is surjective.  Indeed, on $T_pV$,
\[
 d\log|h_a|_p
   =d(\log|P|+|q|^2)_p+\operatorname{Re}(a\cdot dq).
\]
For any $\lambda\in N_p^*S$, extend
$\lambda-d(\log|P|+|q|^2)_p$ to an ambient real linear functional
$\operatorname{Re}(a\cdot dq)$.  The resulting total differential is
$\lambda$, so its restriction to $T_pS$ vanishes and $a$ belongs to
the critical fiber.  Since
$\dim_{\R}N_p^*S=k$, the inverse image of any proper linear
subspace of $N_p^*S$ has measure zero in the critical fiber.
The positive-measure set $\mathcal A_p$ cannot be contained in such
an inverse image.  Selecting successively outside the span already
chosen gives $k$ independent conormals and corresponding functions
$h_1,\ldots,h_k$.

Finally choose the local branches of $\log h_j$ with value zero
at $p$.  Their real differentials vanish on
$\C^{n-k}\times\R^k=T_pS$.  Complex linearity forces the
$dz$ coefficients to be zero and the $dw$ coefficients to be
purely imaginary.  Independence of the real conormals is exactly
independence of the resulting real rows $\beta_j$.
\end{proof}

Fix the resulting $p$, its chart $(z_0,w_0)$, the data
$s_*,\psi,F$ of Proposition~\ref{prop:boundary-data}, and the
support functions for the rest of the proof.
